\uuid{E7O7}
\titre{Résolution d'équation différentielle du premier ordre}
\auteur{bodin ; croizat ; sacre}
\organisation{exo7}
\licence{CC BY-NC-SA 4.0}
\datecreate{2022-03-08}
\niveau{L1}
\module{Analyse}
\chapitre{Equation différentielle}
\sousChapitre{Résolution d'équation différentielle du premier ordre}
\difficulte{2}
\theme{}
\exo7qcmid{hilisit.equations-differentielles.044}

\contenu{
\begin{question}
\qclassification{225.01}
Parmi les fonctions suivantes, quelles sont celles qui sont solutions sur \(\Rr\) de l'équation différentielle \(y'=2y+ e^{2x}\) et qui valent \(2\) en \(x=0\) :
\begin{answers}
	\bad{\(x\mapsto 2 e^{2x}\)}
	\bad{\(x\mapsto x e^{2x}\)}
	\bad{\(x\mapsto x e^{2x}+2\)}
    \good{\(x\mapsto (x+2) e^{2x}\)}
\end{answers}
\begin{explanations}
On peut éliminer la fonction \(x\mapsto x e^{2x}\) qui ne prend pas la valeur \(2\) en \(x=0\) contrairement aux trois autres. On calcule ensuite la dérivée des autres fonctions proposées et on remplace \(y\)  et \(y'\) dans l'équation différentielle pour identifier celle qui est solution : la seule qui soit solution de notre équation différentielle est $x \mapsto (x+2)e^{2x}$. Rappel : la dérivée du produit de deux fonctions \(u\) et \(v\) est \(u'v+uv'\). Ainsi $[(x+2)e^{2x}]' =  e^{2x} + 2(x+2)e^{2x}$.
\end{explanations}
\end{question}
}
